% mazzi: 11,12
\chapter[Lo stimolo vibratorio]{Lo stimolo vibratorio: dalla prestazione alla valutazione del
giocatore infortunato}

\textbf{Stimolo vibratorio}: metodica di stimolazione meccanica del sistema biologico, diffusa in
Italia dalla fine degli anni Novanta, che interessa la valutazione funzionale sotto un duplice
profilo:
\begin{itemize}
  \item come \textbf{mezzo di allenamento}, per gli effetti positivi sulla prestazione e in
    particolare sul sistema neuromuscolare;
  \item come possibile \textbf{strumento di valutazione}, soprattutto nel recupero del soggetto
    infortunato.
\end{itemize}
L'interesse per la metodica è in crescita costante: le pubblicazioni indicizzate su \textit{PubMed}
passano da poche decine nei quinquenni tra il 1960 e il 1990 a diverse centinaia nel periodo
2010--2014, e si estendono dall'ambito prestativo a quello clinico, per gli adattamenti positivi del
sistema neuromuscolare.

% ---------------------------------------------------------------------------
\section{Gli effetti sul sistema biologico}

Effetti osservati in letteratura, sia \textbf{acuti} sia \textbf{cronici}:
\begin{enumerate}
  \item miglioramento della curva \textbf{forza-velocità};
  \item miglioramento della curva \textbf{potenza-velocità};
  \item \textbf{adattamenti neuromuscolari} significativi;
  \item miglioramento della \textbf{forza esplosiva};
  \item miglioramento della \textbf{resistenza alla forza veloce};
  \item aumento della \textbf{flessibilità muscolare} e della \textbf{mobilità articolare};
  \item miglioramento dell'\textbf{equilibrio posturale};
  \item stimolazione del \textbf{sistema endocrino};
  \item stimolazione del \textbf{sistema metabolico}.
\end{enumerate}

\rubrica{Il paradosso forza-flessibilità}
Forza e flessibilità di norma non si muovono nella stessa direzione:
\begin{itemize}
  \item allenamento di forza $\rightarrow$ riduzione della flessibilità;
  \item allenamento della flessibilità $\rightarrow$ diminuzione della forza, almeno in acuto.
\end{itemize}
Con lo stimolo vibratorio i due fenomeni si osservano invece quasi contemporaneamente, già in acuto:
è il dato che richiede una spiegazione meccanicistica (vedi §\ref{sec:perche}).

% ---------------------------------------------------------------------------
\section{Cenni storici}

\rubrica{L'origine ergonomica}
I primi studi nascono nell'\textbf{ergonomia industriale}, con l'obiettivo opposto a quello
sportivo: ridurre il propagarsi della vibrazione nel sistema biologico.
\begin{itemize}
  \item ambiti indagati: carrozze, trapani, martelli pneumatici, comfort dei mezzi pesanti,
    automobili, elicotteri;
  \item il problema è l'esposizione sistematica e prolungata del lavoratore (motoseghe, martelli
    pneumatici, piloti di elicottero) $\rightarrow$ problematiche di tipo neurale.
\end{itemize}
Il sistema biologico riconosce comunque lo stimolo vibratorio perché lo subisce da sempre: nella
corsa la vibrazione attraversa tutto il corpo, e lo stesso avviene in automobile o in moto.

\rubrica{Dalla neurologia allo sport}
\begin{enumerate}
  \item \textbf{anni Sessanta}: l'argomento è appannaggio dei neurologi, che individuano i primi
    comportamenti significativi in seguito a questo tipo di stimolazione;
  \item \textbf{Nazarov} (USSR), anni Ottanta: applica la \textbf{vibrazione locale} all'allenamento
    delle ginnaste (forza e flessibilità), al balletto del Bol'šoj (flessibilità e trattamento degli
    infortuni) e ad applicazioni terapeutiche;
  \item \textbf{Nazarov e Spivak, 1985}: il lavoro viene inizialmente ignorato per ragioni di
    barriera linguistica e politica --- la metodica diventa pubblica solo dopo la caduta del
    comunismo;
  \item \textbf{Issurin}: introduce la \textit{superimposed vibration}, cioè la vibrazione
    sovrapposta all'esercizio;
  \item \textbf{Bosco}: introduce la \textit{whole body vibration} (WBV) e apre la stagione di
    ricerca che sdogana la metodica nell'ambito sportivo.
\end{enumerate}

% ---------------------------------------------------------------------------
\section{Le evidenze sperimentali}

\subsection{Issurin e coll., 1998 --- effetti acuti e residui sulla forza esplosiva}

\rubrica{Metodo}
Vibrazione locale sovrapposta all'esercizio: il dispositivo vibratorio è posto sul cavo tenuto dal
soggetto, cavo a sua volta collegato al carico attraverso un sistema di pulegge; l'esercizio è un
\textit{curl} bilaterale dei bicipiti, strumentato con \textit{power monitor} e sonde.

\rubrica{Risultati}
Confronto tra condizione \textbf{VS} (con stimolo vibratorio) e \textbf{No VS} (controllo), in
atleti elite e amatori, in termini di \textit{power gain}:
\begin{itemize}
  \item \textbf{elite}: guadagno di potenza $\approx +30$\,W con VS, contro valori negativi
    ($\approx -2/-3$\,W) senza VS;
  \item \textbf{amatori}: $\approx +20/+25$\,W con VS, contro $\approx -8$\,W senza VS;
  \item il guadagno di potenza migliora dunque in entrambi i gruppi sottoposti a vibrazione, con
    differenza statisticamente significativa rispetto ai controlli.
\end{itemize}

\subsection{Bosco e coll., 1998 --- WBV e salto verticale}

\rubrica{Metodo}
Piattaforma a \textbf{oscillazione alternata}, esposizione di 8 minuti.

\rubrica{Risultati}
Nel gruppo sperimentale, rispetto al gruppo di controllo:
\begin{itemize}
  \item aumento dell'\textbf{innalzamento del centro di gravità} nel protocollo di salto continuo di
    5 secondi ($p<0{,}005$);
  \item aumento della \textbf{potenza} sviluppata durante il salto, in W/kg ($p<0{,}05$);
  \item aumento del \textbf{miglior salto} registrato nei 5 secondi di salto continuo ($p<0{,}05$);
  \item nel gruppo di controllo nessuna variazione.
\end{itemize}

\rubrica{Evoluzione della macchina}
L'oscillazione alternata di questa prima piattaforma poteva creare problemi a livello pelvico e
della colonna $\rightarrow$ la concezione è stata modificata introducendo l'\textbf{oscillazione
sinusoidale verticale}, tuttora in uso.

\subsection{Bakhtiary e coll., 2007 --- vibrazione e DOMS}

\textbf{DOMS} (\textit{delayed onset muscle soreness}): dolore muscolare ritardato, tipico
dell'inizio della stagione sportiva.
\begin{itemize}
  \item insorge in seguito a lavoro di tipo \textbf{eccentrico} e dura da 24--36 fino a un massimo
    di 48 ore;
  \item è dovuto alle \textbf{microlesioni} a livello del tessuto muscolare, di per sé positive
    $\rightarrow$ \textit{sarcomerogenesi}, cioè l'adattamento al lavoro iniziale.
\end{itemize}

\rubrica{Metodo}
Cinquanta volontari sani non atleti, assegnati casualmente a due gruppi: \textbf{VT} ($n=25$) e
\textbf{non-VT} ($n=25$).
\begin{enumerate}
  \item nel gruppo VT si applica una vibrazione a 50\,Hz per 1 minuto su quadricipiti, ischiocrurali
    e polpacci, a destra e a sinistra; nel gruppo non-VT nessuna vibrazione;
  \item entrambi i gruppi camminano poi in \textbf{discesa} su treadmill inclinato di $10^\circ$
    alla velocità di 4\,km/h.
\end{enumerate}

\rubrica{Risultati}
\begin{itemize}
  \item \textbf{forza} maggiore nel gruppo VT, sia nell'arto destro sia nel sinistro;
  \item \textbf{percezione del dolore} ridotta nel gruppo VT sulla scala utilizzata;
  \item soprattutto, \textbf{creatinchinasi} (CK) plasmatica --- enzima che compare nel sangue in
    seguito a microlesione muscolare --- nettamente più bassa nel gruppo VT: $\approx 115$ contro
    $\approx 195$\,IU/l ($p=0{,}001$).
\end{itemize}

\subsection{Bosco e coll., 2001 --- due mesi di WBV su calciatori}

Adattamenti neuromuscolari in calciatori dopo 2 mesi di WBV, con test ripetuto a luglio e ad agosto.

\begin{table}[htbp]
  \centering
  \small
  \renewcommand{\arraystretch}{1.3}
  \begin{tabularx}{\textwidth}{|X|c|c|}
    \hline
    \textbf{Test} & \textbf{T. luglio (cm)} & \textbf{T. agosto (cm)} \\
    \hline
    CMJ (salto con contromovimento) & $\approx 44$ & $\approx 46$ \\
    \hline
    CMJL (con l'uso delle braccia) & $\approx 51{,}5$ & $\approx 52$ \\
    \hline
    5 CMJ (5 secondi di salto continuo) & $\approx 35$ & $\approx 39$ \\
    \hline
    15\,$''$ (resistenza alla forza veloce) & $\approx 35$ & $\approx 39$ \\
    \hline
    \textit{Sit and reach} (flessibilità) & $\approx 39{,}5$ & $\approx 50$ \\
    \hline
  \end{tabularx}
\end{table}

Il dato rilevante è che al miglioramento delle espressioni di forza si accompagna, per la prima
volta, un concomitante miglioramento della \textbf{flessibilità muscolare}.

\subsection{Bosco e coll., 2000 --- le risposte ormonali}

Variazioni della prestazione neuromuscolare e dei livelli ormonali plasmatici in 14 giovani atleti
maschi.

\rubrica{Protocollo}
60 secondi di WBV seguiti da 60 secondi di recupero, ripetuti 10 volte (10 minuti complessivi di
vibrazione).

\rubrica{Risposte osservate (effetto acuto)}
\begin{itemize}
  \item \textbf{testosterone}: aumentato;
  \item \textbf{ormone della crescita} (GH): aumentato;
  \item \textbf{cortisolo}: diminuito --- è l'ormone dello stress, normalmente molto elevato a fine
    allenamento e rintracciabile fino a 24--36 ore dopo, espressione della condizione catabolica
    dell'atleta;
  \item \textbf{rapporto T/C}: aumentato in modo significativo, come conseguenza diretta dei due
    andamenti opposti.
\end{itemize}

\textbf{Rapporto testosterone/cortisolo}: rapporto tra l'ormone anabolico e quello catabolico; è
l'indice usato per capire quale dei due aspetti prevalga nello stato biologico del soggetto.

\subsection{Ritzmann e coll., 2014 --- controllo posturale e resistenza muscolare}

Quaranta atleti allenati, 4 settimane di WBV, 3 volte a settimana, 2 serie per sessione, recupero
tra le serie di 3 minuti.

Rispetto al gruppo di controllo non sottoposto a vibrazione si osserva un miglioramento del
\textbf{controllo posturale} (riduzione dell'area di oscillazione del centro di pressione nei piani
antero-posteriore e medio-laterale) e della \textbf{resistenza alla forza}. La vibrazione non
migliora dunque soltanto forza esplosiva e flessibilità, ma anche l'equilibrio posturale.

% ---------------------------------------------------------------------------
\section{Perché migliorano forza esplosiva e flessibilità?}
\label{sec:perche}

Il problema posto dalla letteratura: la vibrazione migliora \textbf{forza esplosiva} e
\textbf{flessibilità}, ma non produce cambiamenti sulla \textbf{forza isometrica massima}
($FI_{max}$), o comunque l'effetto sulla forza massima è ridotto.

\subsection{Il riflesso tonico da vibrazione}

\textbf{Riflesso tonico da vibrazione} (\textit{tonic vibration reflex}): iperattivazione riflessa
dei \textbf{muscoli agonisti} durante l'esposizione allo stimolo.
\begin{itemize}
  \item il soggetto non si muove: mantiene la posizione per tutta la durata dello stimolo;
  \item l'attività EMG (RMS) si innalza nettamente all'inizio della vibrazione, si mantiene per
    tutti i 60 secondi e torna ai valori di base alla fine dello stimolo;
  \item è generato da un'\textbf{attivazione riflessa}, non volontaria.
\end{itemize}

\subsection{Annino e coll. --- che cosa succede dopo la vibrazione}

\rubrica{Metodo}
\begin{itemize}
  \item 20 giovani adulti maschi ($24{,}8 \pm 2{,}5$ anni), randomizzati in due gruppi;
  \item \textbf{VG} (gruppo vibrazione): 10 minuti di WBV a 35\,Hz, spostamento 5\,mm, magnitudo
    5\,g;
  \item \textbf{NVG} (gruppo non vibrato, placebo): stessa posizione sulla pedana, senza alcuna
    vibrazione;
  \item valutazione con \textbf{CMJ} e flessibilità, con analisi EMG su vasto laterale (VL), vasto
    mediale (VM), bicipite femorale (BF) e gastrocnemio laterale (LG).
\end{itemize}

\rubrica{Risultati}
Confronto tra il protocollo breve (5$'$) e quello lungo (5+5$'$, cioè 10 minuti):

\begin{table}[htbp]
  \centering
  \small
  \renewcommand{\arraystretch}{1.3}
  \begin{tabularx}{\textwidth}{|X|c|c|c|c|}
    \hline
    \multirow{2}{*}{\textbf{Variabile}} & \multicolumn{2}{c|}{\textbf{Prova di salto}} &
      \multicolumn{2}{c|}{\textbf{Prova di flessibilità}} \\
    \cline{2-5}
     & \textbf{5$'$} & \textbf{5+5$'$} & \textbf{5$'$} & \textbf{5+5$'$} \\
    \hline
    Quadricipite (EMG) & invariato & invariato & invariato & invariato \\
    \hline
    Ischiocrurali (EMG) & $\downarrow$ ** & $\downarrow$ ** & $\downarrow$ * & $\downarrow$ ** \\
    \hline
    Rapporto H/Q & $\downarrow$ * & $\downarrow$ ** & --- & --- \\
    \hline
    Gastrocnemio (EMG) & invariato & invariato & invariato & invariato \\
    \hline
    Altezza del CMJ & invariata & $\uparrow$ ** & --- & --- \\
    \hline
    Flessibilità (spostamento) & --- & --- & $\uparrow$ * & $\uparrow$ ** \\
    \hline
  \end{tabularx}
\end{table}

\rubrica{Interpretazione}
\begin{enumerate}
  \item il miglioramento del salto non è dovuto a una maggiore attivazione del
    \textbf{quadricipite}: subito dopo la vibrazione l'innalzamento dell'agonista non è
    statisticamente significativo, contrariamente a quanto si riteneva osservando il riflesso tonico
    da vibrazione durante lo stimolo;
  \item si osserva invece una riduzione dell'attività dei flessori, cioè degli \textbf{antagonisti};
  \item riduzione degli antagonisti $\rightarrow$ miglior \textbf{rapporto H/Q} $\rightarrow$
    maggior altezza di salto;
  \item l'effetto sul salto non compare con 5 minuti di vibrazione, ma compare con 10 minuti;
  \item anche l'aumento della flessibilità, in entrambi i protocolli, è dovuto alla riduzione
    dell'attivazione dei flessori più che a un miglioramento dell'azione degli agonisti.
\end{enumerate}
\textbf{Conclusione}: la vibrazione agisce sul comportamento neuromuscolare attraverso gli
\textbf{effetti inibitori} sui muscoli antagonisti, più che attraverso l'attività riflessa da
stiramento sugli agonisti. Nel gruppo di controllo nessuna differenza significativa in tutti i
parametri considerati.

\rubrica{Effetti sul ciclo stiramento-accorciamento}
Nell'analisi del CMJ si osserva inoltre una diminuzione del \textbf{tempo eccentrico} ($T_{ecc}$) e
del \textbf{tempo concentrico} ($T_{conc}$), cioè una riduzione della durata complessiva del
\textbf{ciclo stiramento-accorciamento}. La riduzione è marcata nel protocollo lungo ($\approx
-13\%$ e $\approx -15\%$), dove si accompagna all'unico incremento dell'altezza di salto ($\approx
+4\%$).

Il principio generale è che nelle prestazioni non conta solo l'azione degli agonisti, ma il
\textbf{rapporto tra agonisti e antagonisti}, che interviene in modo determinante anche nei
movimenti esplosivi.

% ---------------------------------------------------------------------------
\section{Lo stimolo vibratorio associato all'esercizio fisico}

Seconda modalità applicativa: non la vibrazione isolata, ma la vibrazione \textbf{sovrapposta
all'esercizio}.

\subsection{Ronnestad, 2004 --- WBV e squat training}

Confronto tra due gruppi che lavorano con i sovraccarichi, uno con vibrazione e uno senza;
allenamento 3 volte a settimana per 5 settimane.

\begin{table}[htbp]
  \centering
  \small
  \renewcommand{\arraystretch}{1.3}
  \begin{tabularx}{\textwidth}{|X|c|c|c|c|}
    \hline
    \multirow{2}{*}{\textbf{Misura}} & \multicolumn{2}{c|}{\textbf{WBV \& squat training}} &
      \multicolumn{2}{c|}{\textbf{Squat training}} \\
    \cline{2-5}
     & \textbf{pre} & \textbf{post} & \textbf{pre} & \textbf{post} \\
    \hline
    1RM --- forza (N) & $\approx 160$ & $\approx 215$ & $\approx 148$ & $\approx 182$ \\
    \hline
    CMJ --- innalzamento del C.G. (cm) & $\approx 36$ & $\approx 39{,}5$ & $\approx 34{,}5$ &
      $\approx 36$ \\
    \hline
  \end{tabularx}
\end{table}

\begin{itemize}
  \item sulla \textbf{forza dinamica massima} (1RM) il miglioramento è maggiore nel gruppo che
    associa vibrazione e squat;
  \item sul salto la differenza è ancora più netta: il miglioramento del gruppo con vibrazione è di
    gran lunga superiore a quello del gruppo che ha svolto il solo squat training.
\end{itemize}

\subsection{Annino e coll., 2018 --- WBV e \textit{repeated sprint ability}}

\begin{itemize}
  \item rispetto al protocollo classico, che prevede una sola serie di vibrazione, sono state
    utilizzate 3 serie, così da prolungare l'effetto e determinare una fatica aggiuntiva;
  \item nel gruppo sottoposto a vibrazione (SG) il \textbf{tempo medio della RSA} cresce molto meno
    che nel gruppo di controllo (CG) lungo tutte e tre le serie: da $\approx 7{,}95$ a $\approx
    8{,}1$\,s contro $\approx 8{,}0$ $\rightarrow$ $\approx 8{,}4$\,s;
  \item il gruppo vibrato mantiene più a lungo la prestazione su un metabolismo di tipo anaerobico,
    coinvolgendo le \textbf{unità motorie veloci} per molto più tempo: la vibrazione le condiziona
    quindi in misura maggiore.
\end{itemize}

% ---------------------------------------------------------------------------
\section{Come agisce la vibrazione}

\subsection{Il tempo di esposizione}

\textbf{Allenamento}: programmazione di uno stimolo meccanico generato attraverso modificazioni
delle condizioni ambientali (sovraccarico, pliometria, \dots). Il presupposto dell'adattamento del
sistema biologico allo stimolo meccanico è il \textbf{tempo di esposizione}.

\rubrica{Il confronto pliometria / sovraccarichi}
\begin{itemize}
  \item nessuno cerca l'\textbf{ipertrofia} con la pliometria, pur essendo questa estremamente
    intensa: il tempo di lavoro muscolare è troppo breve;
  \item il lavoro con i sovraccarichi raggiunge invece tempi di contrazione di 700--900\,ms e oltre
    --- fino a $1$--$1{,}2$\,s con carichi molto alti --- e determina per questo un adattamento di
    tipo \textbf{morfologico};
  \item nella progressione degli adattamenti i tempi di contrazione scalano dalla forza massima
    ($700$--$900$\,ms) alla forza dinamica massima ($200$--$400$\,ms), alla forza esplosiva
    ($150$--$200$\,ms), fino alla forza esplosiva prestativa ($\approx 100$\,ms);
  \item nella corsa la \textbf{forza di reazione verticale del terreno} cresce con la velocità ---
    da $\approx 1000$\,N a 2\,m/s fino a oltre $3000$\,N a 12\,m/s --- mentre il tempo di contatto
    si accorcia progressivamente.
\end{itemize}

\rubrica{Il vantaggio specifico della vibrazione}
Confrontando l'\textbf{accelerazione media} (in G) con il \textbf{tempo di applicazione} (in ms)
nelle diverse prestazioni:
\begin{enumerate}
  \item i 5--6\,G generati dalla macchina vibrante si ritrovano, nello sport, solo in prestazioni
    come il \textbf{salto in lungo} ($\approx 6$\,G) e il \textbf{salto in alto} ($\approx 5$\,G),
    al momento dell'impatto al suolo;
  \item in quei gesti però il tempo di applicazione è di circa 100\,ms; i salti con sovraccarico,
    che durano molto di più (fino a $800$--$900$\,ms), esprimono accelerazioni di appena
    $1$--$1{,}5$\,G;
  \item con la vibrazione si può invece sottoporre il sistema biologico a 5--6\,G per un tempo molto
    più lungo: è la combinazione altrimenti irraggiungibile di magnitudo elevata e lunga
    esposizione.
\end{enumerate}

\rubrica{Perché non danneggia le articolazioni}
Un tempo equivalente di lavoro pliometrico produrrebbe problemi articolari e muscolari gravissimi.
La differenza sta nell'\textbf{ampiezza dello spostamento}:
\begin{itemize}
  \item nella piattaforma vibrante l'oscillazione riguarda pochi millimetri, e lo stress
    sull'articolazione è quindi minimo;
  \item nel lavoro pliometrico l'\textbf{allungamento del tendine} può raggiungere i 2--3
    centimetri.
\end{itemize}

\subsection{I parametri del protocollo}

\rubrica{Linee guida generali}
\begin{itemize}
  \item \textbf{durata della vibrazione}: 10--60 secondi per singola ripetizione;
  \item \textbf{recupero}: sempre 1 minuto tra una ripetizione e l'altra;
  \item \textbf{volume}: massimo 10 serie;
  \item \textbf{frequenza}: massimo 1 sessione al giorno.
\end{itemize}

\rubrica{Parametri della vibrazione}
\begin{enumerate}
  \item \textbf{ampiezza}: 4--6\,mm;
  \item \textbf{frequenza}: 20--55\,Hz;
  \item \textbf{magnitudo}: $> 5$\,g.
\end{enumerate}

\rubrica{Altri parametri da controllare}
\begin{itemize}
  \item \textbf{esposizione totale} al trattamento;
  \item \textbf{posizione del corpo} sulla pedana.
\end{itemize}

% ---------------------------------------------------------------------------
\section{L'integrazione sensomotoria come substrato}

Il riflesso tonico da vibrazione si spiega con l'\textbf{integrazione sensomotoria}: l'iperattività
osservata è verosimilmente dovuta alla stimolazione dei \textbf{meccanocettori} e
\textbf{propriocettori} a livello muscolare, tendineo e articolare, che generano un \textbf{arco
riflesso} e aumentano così l'attività riflessa e l'azione muscolare a livello spinale.

\rubrica{I recettori muscolari}
\begin{itemize}
  \item \textbf{fusi neuromuscolari};
  \item \textbf{organi tendinei del Golgi} (GTO);
  \item \textbf{corpuscoli paciniformi};
  \item \textbf{terminazioni nervose libere} (FNE).
\end{itemize}

\rubrica{I recettori articolari e legamentosi}
La funzione muscolare è integrata anche dalle informazioni afferenti provenienti dai
\textbf{recettori articolari} --- corpuscoli di Pacini, di Ruffini, organi \textit{Golgi-like} e
terminazioni nervose libere. Scoperti intorno agli anni Novanta, sono presenti anche nei legamenti
delle articolazioni.

Dalle afferenze legamentose i \textbf{legamenti} agiscono per due vie distinte:
\begin{itemize}
  \item per le loro \textbf{proprietà meccaniche} $\rightarrow$ direttamente sulla \textbf{stabilità
    articolare};
  \item per le loro \textbf{proprietà sensoriali} $\rightarrow$ su due percorsi:
    \begin{itemize}
      \item sul \textbf{sistema $\gamma$--fuso neuromuscolare} $\rightarrow$ controllo della
        \textit{stiffness} muscolare e della coordinazione $\rightarrow$ stabilità articolare;
      \item sul \textbf{senso di posizione e di movimento}.
    \end{itemize}
\end{itemize}

I recettori legamentosi forniscono dunque il senso di posizione e la stabilità articolare, oltre a
controllare la coordinazione e la \textit{stiffness} muscolare attraverso il sistema $\gamma$ dei
fusi neuromuscolari.

\rubrica{Il parallelo con il riflesso da stiramento}
Il potenziamento dell'attivazione dei nervi motori passa per il \textbf{riflesso miotatico}
(riflesso da stiramento).
\begin{enumerate}
  \item lo stimolo meccanico vibratorio si comporta come il riflesso da stiramento;
  \item porta all'attivazione dell'agonista attraverso la stimolazione del \textbf{fuso
    neuromuscolare};
  \item con la concomitante \textbf{inibizione} dei muscoli antagonisti $\rightarrow$ è la
    spiegazione del miglioramento simultaneo di forza e flessibilità;
  \item lo stimolo raggiunge inoltre livelli nervosi più alti, fino a stimolare l'\textbf{ipotalamo}
    $\rightarrow$ variazioni ormonali (GH) che contribuiscono a loro volta all'aumento della forza
    muscolare.
\end{enumerate}

\rubrica{Prospettiva applicativa}
La ricchezza di informazioni che lo stimolo vibratorio fornisce sul sistema riflesso lo rende un
candidato come \textbf{sistema di valutazione} nel recupero, soprattutto nei soggetti infortunati: è
l'applicazione sviluppata nel resto del capitolo.

% ---------------------------------------------------------------------------
\section{Il problema del giocatore infortunato}

\rubrica{Il limite del parametro biomeccanico}
La prestazione non è espressione dei soli \textbf{parametri biomeccanici}: coinvolge il sistema
nervoso in tutto il suo complesso.
\begin{itemize}
  \item dietro i parametri misurabili sta una serie di eventi generata a livello delle \textbf{aree
    motorie cerebrali} --- l'area motoria 4 in interazione con le altre aree --- e integrata con
    tutte le afferenze provenienti dal resto del corpo;
  \item il sistema nervoso governa il movimento attraverso quattro sottosistemi: \textbf{sistema
    propriocettivo}, \textbf{sistema esterocettivo}, \textbf{sistema piramidale} e \textbf{sistema
    extrapiramidale};
  \item un infortunio non riguarda quindi il solo aspetto biomeccanico: crea \textbf{adattamenti a
    livello neurale} che compromettono l'espressione del sistema neuromuscolare.
\end{itemize}
Ne segue che stabilire se un giocatore ha recuperato non può basarsi solo su un parametro
biomeccanico misurato in un test semplice: nella prestazione il sistema neuromuscolare dovrà
eseguire movimenti molto più complessi di quelli del test.

\rubrica{Il loop dell'infortunio}
Il giocatore infortunato entra in un circuito da cui esce con difficoltà, perché l'infortunio stesso
è il \textbf{principale fattore di rischio} di un infortunio successivo --- recidive, o infortuni in
altre sedi corporee per effetto di \textbf{atteggiamenti compensatori}. Il percorso, a partire
dall'impatto, si articola così:
\begin{enumerate}
  \item \textbf{infortunio} $\rightarrow$ eventuale \textbf{intervento chirurgico} $\rightarrow$
    \textbf{riabilitazione}; oppure direttamente riabilitazione;
  \item la riabilitazione da sola non restituisce il recupero funzionale: l'\textbf{acquisizione
    clinica} non coincide quasi mai con quella funzionale;
  \item serve una fase di \textbf{riatletizzazione}, in cui l'atleta riacquisisce le abilità motorie
    e tecnico-sportive ai ritmi e alle intensità di gara, condizionando il sistema neuromuscolare ad
    azioni di altissima intensità e complessità;
  \item solo allora si torna alla prestazione. Se la riatletizzazione non avviene correttamente, si
    rientra nel loop.
\end{enumerate}

\subsection{Flessori ed estensori del ginocchio: il comportamento per angolo}

Uno degli aspetti da considerare è il diverso comportamento meccanico di flessori ed estensori
dell'arto inferiore al variare dell'angolo articolare. Gli angoli bassi corrispondono alla gamba
quasi distesa, i $90^\circ$ alla massima flessione.

\begin{table}[htbp]
  \centering
  \small
  \renewcommand{\arraystretch}{1.3}
  \begin{tabularx}{\textwidth}{|X|c|c|c|c|c|c|c|}
    \hline
    \textbf{Angolo} & $5^\circ$ & $15^\circ$ & $30^\circ$ & $45^\circ$ & $60^\circ$ & $75^\circ$ &
      $90^\circ$ \\
    \hline
    Momento flessori (Nm) & 63,6 & 57,4 & 56,9 & 49,5 & 50,5 & 45,7 & 36,1 \\
    \hline
    Momento estensori (Nm) & 61,5 & 85,5 & 107,4 & 120,9 & 119,5 & 117 & 103,9 \\
    \hline
    Forza flessori (N) & 2072 & 1772 & 1515 & 1238 & 1299 & 1322 & 1506 \\
    \hline
    Forza estensori (N) & 1418 & 1814 & 2213 & 2479 & 2576 & 2727 & 2768 \\
    \hline
  \end{tabularx}
\end{table}

La forza è ricavata dal momento secondo $F = M / l_m$, dove $l_m$ è il braccio di leva.

\rubrica{Lettura}
\begin{itemize}
  \item l'\textbf{estensore} (quadricipite) raggiunge il picco di momento intorno ai $45^\circ$ e da
    lì decresce;
  \item i \textbf{flessori} (ischiocrurali: bicipite femorale, gracile, semimembranoso,
    semitendinoso) hanno il picco a gamba distesa, non verso la fine del movimento;
  \item l'andamento è lo stesso considerando la forza al posto del momento: l'azione dei flessori è
    molto più elevata a gamba quasi distesa.
\end{itemize}

\rubrica{Conseguenza allenante}
Gli ischiocrurali intervengono quasi all'inizio del movimento --- in alcuni gesti esplosivi e nella
corsa entrano quando la gamba è in massima estensione --- cioè nella loro fase di
\textbf{allungamento} e non di accorciamento. Vanno perciò allenati in \textbf{eccentrico} e non in
concentrico, come invece si fa abitualmente: il \textit{Nordic hamstring curl} è per questo tra gli
esercizi da preferire.

% ---------------------------------------------------------------------------
\section{L'epidemiologia degli infortuni nel calcio}

\subsection{La distribuzione temporale durante la gara}

\rubrica{Il dato storico (Woods e coll., 2003)}
Distribuzione percentuale dei traumi distorsivi della \textbf{tibio-tarsica} secondo la collocazione
durante la gara: l'incidenza cresce verso la fine di entrambi i tempi --- $\approx 24\%$ nel quarto
d'ora 31$'$--45$'$ e $\approx 24\%$ nel 76$'$--90$'$, contro l'$8\%$ del primo quarto d'ora ---
andamento riconducibile a una condizione di \textbf{fatica}.

\rubrica{Il dato recente}
Un lavoro più recente del gruppo del docente, sui quattro maggiori campionati europei, mostra un
quadro diverso:
\begin{itemize}
  \item nel \textbf{primo tempo} il comportamento è simile ovunque: incidenza massima verso la fine,
    nel quarto d'ora 31$'$--45$'$;
  \item nel \textbf{secondo tempo} il picco si sposta a seconda del campionato --- nel 61$'$--75$'$
    per Serie A e Premier League, negli ultimi 15 minuti per Liga e Bundesliga;
  \item complessivamente il secondo tempo risulta meno coinvolto di quanto indicasse il lavoro del
    2003.
\end{itemize}

\subsection{Cinque stagioni di Serie A}

Tipi di infortunio nelle ultime cinque stagioni di Serie A:

\begin{table}[htbp]
  \centering
  \small
  \renewcommand{\arraystretch}{1.3}
  \begin{tabularx}{\textwidth}{|X|c|c|c|c|c|}
    \hline
    \textbf{Tipo} & \textbf{2012/13} & \textbf{2013/14} & \textbf{2014/15} & \textbf{2015/16} &
      \textbf{2016/17} \\
    \hline
    Da contatto & 132 & 127 & 113 & 114 & 112 \\
    \hline
    Non da contatto & 349 & 317 & 288 & 355 & 338 \\
    \hline
    Altro & 84 & 96 & 72 & 73 & 50 \\
    \hline
    LCA & 12 & 19 & 16 & 17 & 18 \\
    \hline
  \end{tabularx}
\end{table}

\begin{itemize}
  \item in tutte e cinque le stagioni l'infortunio \textbf{non da contatto} è di gran lunga il più
    frequente;
  \item il \textbf{legamento crociato anteriore} (LCA) resta pressoché stabile, con una media di
    16--17 casi per stagione.
\end{itemize}

\rubrica{Perché il dato non da contatto è il più interessante}
Un muscolo, un tendine o un'articolazione nell'esercizio delle proprie funzioni non dovrebbero
rompersi: se accade senza un contatto esterno c'è un'\textbf{alterazione} del comportamento del
sistema neuromuscolare, incompatibile con la prestazione ad altissima velocità. È il presupposto che
giustifica una valutazione di tipo neuromuscolare, e non solo biomeccanica.

\subsection{Il confronto tra i quattro maggiori campionati europei}

\rubrica{Contatto e non contatto}
Il predominio del non contatto si conferma ovunque, ma la percentuale italiana è la più alta:
\begin{itemize}
  \item \textbf{Serie A}: 32\% da contatto, 68\% non da contatto;
  \item \textbf{Premier League}: 36\% / 64\%;
  \item \textbf{Liga}: 38\% / 62\%;
  \item \textbf{Bundesliga}: 37\% / 63\%.
\end{itemize}

\rubrica{Infortuni per ruolo}
I \textbf{difensori} sono il ruolo più colpito, con una distribuzione interna tra centrali ed
esterni; \textbf{centrocampisti} e \textbf{attaccanti} presentano un'incidenza minore.

\rubrica{Tipologia: muscolare o articolare}
\begin{table}[htbp]
  \centering
  \small
  \renewcommand{\arraystretch}{1.3}
  \begin{tabularx}{\textwidth}{|X|c|c|c|c|}
    \hline
    \textbf{Tipologia} & \textbf{Serie A} & \textbf{Premier League} & \textbf{Liga} &
      \textbf{Bundesliga} \\
    \hline
    Muscolare & 104 & 88 & 65 & 58 \\
    \hline
    Articolare & 21 & 32 & 29 & 19 \\
    \hline
  \end{tabularx}
\end{table}

Gli infortuni muscolari superano ovunque gli articolari, ma lo squilibrio è massimo in Serie A.

\rubrica{Tempi di recupero}
Medie in giorni, con il numero di casi tra parentesi:

\begin{table}[htbp]
  \centering
  \small
  \renewcommand{\arraystretch}{1.3}
  \begin{tabularx}{\textwidth}{|X|c|c|}
    \hline
    \textbf{Campionato} & \textbf{Infortuni gravi ($>30$ gg)} & \textbf{LCA (non da contatto)} \\
    \hline
    Serie A & 71,8 (30 casi) & 149 (6 casi) \\
    \hline
    Premier League & 88,3 (37 casi) & 334,6 (3 casi) \\
    \hline
    Liga & 128,5 (30 casi) & 252,1 (10 casi) \\
    \hline
    Bundesliga & 75,5 (21 casi) & 190,5 (2 casi) \\
    \hline
  \end{tabularx}
\end{table}

\subsection{Il confronto con la letteratura e le due criticità}

Tra gli studi sul tema il più vicino a questo lavoro è lo \textit{UEFA injury study}. Le analogie
riguardano in particolare:
\begin{itemize}
  \item la \textbf{distribuzione temporale} degli infortuni durante le gare;
  \item il \textbf{rapporto tra infortuni muscolari e articolari}.
\end{itemize}

Lo studio delinea andamenti ricorrenti, analogie e differenze tra i «top 4» campionati europei,
tracciandone in parte il \textbf{modello prestativo} pur considerando i soli casi di infortunio e
non dati tecnico-tattici. Due aspetti, in particolare, andrebbero approfonditi con ulteriori studi:
\begin{enumerate}
  \item l'alto numero di \textbf{infortuni muscolari} in Serie A;
  \item la \textbf{precocità del recupero} da infortuni gravi in Serie A --- il dato più basso tra i
    quattro campionati, che dovrebbe far riflettere.
\end{enumerate}

% ---------------------------------------------------------------------------
\section{La valutazione neuromuscolare del giocatore infortunato}

\subsection{Dietro il parametro biomeccanico: l'EMG del salto}

L'analisi di un salto verticale restituisce i parametri biomeccanici --- forza verticale, velocità,
spostamento, altezza di salto e potenza --- scanditi nelle fasi del gesto: \textit{standing},
eccentrica, concentrica, volo, atterraggio, ritorno alla posizione eretta. L'esame elettromiografico
sincronizzato aggiunge però un piano di informazione che i soli parametri meccanici non danno.

\rubrica{Il tracciato}
EMG (RMS) registrato insieme a velocità e posizione del carico su vasto laterale e vasto mediale di
entrambi gli arti, ischiocrurali di entrambi gli arti, gastrocnemio destro.

\rubrica{Che cosa rivela in un soggetto infortunato}
\begin{enumerate}
  \item \textbf{attivazione precoce} dell'ischiocrurale dell'arto infortunato durante la fase
    concentrica, rispetto al controlaterale;
  \item \textbf{preattivazione} del flessore poco prima della ricaduta;
  \item ne risulta un alto livello di \textbf{co-contrazione}: aumento dell'attività del flessore
    $\rightarrow$ riduzione della velocità e della prestazione di quell'arto rispetto al
    controlaterale;
  \item si genera così un'\textbf{asimmetria} tra i due arti che nessun test biomeccanico di salto o
    di spinta riesce a rilevare.
\end{enumerate}
È il motivo per cui l'esame diventa decisivo: permette di capire quali \textbf{strategie
neuromuscolari} il sistema ha adottato dopo l'infortunio.

\subsection{Lo stimolo vibratorio nei soggetti operati di LCA}

Studio del 1999--2000 sulla valutazione delle proprietà neuromuscolari in riabilitazione.

\rubrica{Il risultato di base}
Confronto dell'attività EMG (\% EMG RMS) tra arto \textbf{non operato} (N) e arto \textbf{operato}
(O):
\begin{itemize}
  \item \textbf{prima} della vibrazione i due arti sono equivalenti (differenza non significativa);
  \item \textbf{sotto} vibrazione l'arto operato sviluppa un'\textbf{iperattività} rispetto al
    controlaterale ($p<0{,}001$).
\end{itemize}

\rubrica{L'andamento nel tempo}
La differenza tra condizione basale e condizione vibrata, misurata sul quadricipite a 1, 3 e 6 mesi
dall'intervento, resta sistematicamente più alta nell'arto operato:

\begin{table}[htbp]
  \centering
  \small
  \renewcommand{\arraystretch}{1.3}
  \begin{tabularx}{\textwidth}{|X|c|c|c|}
    \hline
    \textbf{Differenza basale/vibrato (\%)} & \textbf{1 mese} & \textbf{3 mesi} & \textbf{6 mesi} \\
    \hline
    Arto operato & $\approx 57$ & $\approx 54$ & $\approx 46$ \\
    \hline
    Arto non operato & $\approx 43$ & $\approx 39$ & $\approx 34$ \\
    \hline
    Significatività & ** & ** & * \\
    \hline
  \end{tabularx}
\end{table}

Il dato critico è il \textbf{disaccoppiamento} tra i due piani di valutazione: già a 3 mesi questi
soggetti presentano test biomeccanici quasi sovrapponibili tra i due arti --- forza, \textit{range}
di movimento, forza isometrica, forza dinamica --- mentre dal punto di vista dell'attivazione
riflessa il deficit è ancora chiaramente presente, e lo resta a 6 mesi.

\rubrica{La spiegazione}
\begin{enumerate}
  \item tendini e legamenti sono ricchi di meccanocettori;
  \item il legamento asportato $\rightarrow$ l'articolazione perde una fonte di informazione
    afferente;
  \item il sistema neuromuscolare si trova privo di un'attività afferente importante per la
    stabilizzazione;
  \item ne consegue l'iperattività dell'arto operato e un cambiamento dell'attività muscolare.
\end{enumerate}

\rubrica{Il modello di valutazione}
Nella valutazione dell'atleta ricostruito di LCA, tre ingressi --- la \textbf{perturbazione}, i
\textbf{meccanocettori articolari} e lo \textbf{stimolo vibratorio} --- agiscono attraverso il fuso
neuromuscolare, gli organi tendinei del Golgi, i motoneuroni $\alpha$ e $\gamma$ e i centri
superiori, e producono variazioni dell'attività muscolare misurabili con la sEMG, isoinerziale e
sotto vibrazione.

\subsection{La vibrazione come mezzo di recupero: Berschin e coll., 2014}

\rubrica{Disegno}
Sessanta soggetti valutati per l'eleggibilità, 20 esclusi, 40 randomizzati in due gruppi da 20
(\textbf{WBV} e \textbf{ST}, protocollo standard); protocollo di esercizio dalla settimana 2 alla
settimana 11 dopo l'intervento, con valutazioni alla 2\textsuperscript{a}, 5\textsuperscript{a},
8\textsuperscript{a} e 11\textsuperscript{a} settimana.

\rubrica{Indice di stabilità (media $\pm$ DS)}
\begin{table}[htbp]
  \centering
  \small
  \renewcommand{\arraystretch}{1.3}
  \begin{tabularx}{\textwidth}{|X|c|c|c|}
    \hline
    \textbf{Settimana} & \textbf{Programma WBV} & \textbf{Programma standard} & \textbf{$p$} \\
    \hline
    2 & $4{,}7 \pm 2{,}3$ & $5{,}4 \pm 3{,}0$ & 0,17 \\
    \hline
    5 & $4{,}0 \pm 1{,}8$ & $5{,}1 \pm 2{,}4$ & 0,07 \\
    \hline
    8 & $3{,}3 \pm 1{,}5$ & $4{,}9 \pm 2{,}4$ & 0,02 * \\
    \hline
    11 & $3{,}1 \pm 1{,}3$ & $4{,}7 \pm 2{,}8$ & 0,01 * \\
    \hline
  \end{tabularx}
\end{table}

I soggetti operati di LCA e sottoposti a WBV per 11 settimane mostrano valori di efficienza
neuromuscolare superiori a quelli del protocollo standard, con differenza che diventa significativa
dall'ottava settimana. La conclusione degli autori è prudente: i dati preliminari indicano che il
protocollo di esercizio con WBV sembra una buona alternativa al programma standard nella
riabilitazione dell'LCA.

\subsection{Il quadro d'insieme}

Lo stimolo vibratorio si colloca nel percorso dell'infortunio come strumento a \textbf{doppia
funzione}, tra la fase di riabilitazione e quella di \textbf{rieducazione funzionale sportiva} che
riporta alla prestazione, con la \textbf{prevenzione} a chiudere il circuito.
\begin{enumerate}
  \item come sistema di valutazione della funzionalità propriocettiva dei meccanocettori, cioè del
    sistema afferente, attraverso il riflesso tonico da vibrazione: nell'arto operato esso manca
    quasi del tutto, e l'alterazione dell'attività elettromiografica che ne deriva è la stessa che
    ci si ritroverà poi durante la prestazione;
  \item come mezzo di allenamento nella fase post-infortunio, dove una stimolazione adeguata può
    rendere il sistema neuromuscolare più efficiente, tenuto conto del deficit in cui si trova dopo
    l'operazione.
\end{enumerate}
Il punto di fondo: dopo un intervento ciò che manca non è più la tenuta biomeccanica dell'arto ---
la chirurgia ortopedica ha fatto passi da gigante --- ma il \textbf{recupero neuro-funzionale}. È su
quel piano che va osservato il fenomeno, e su quel piano che andrebbero fissati i tempi di
\textbf{ritorno in campo}, che non dovrebbero basarsi solo sul recupero biomeccanico.
